\documentclass[11pt,a4paper]{article}
\usepackage[utf8]{inputenc}
\usepackage[T1]{fontenc}
\usepackage{amsmath,amssymb,amsthm}
\usepackage{mathtools}
\usepackage{physics}
\usepackage[margin=2.5cm]{geometry}
\usepackage{hyperref}
\usepackage{cleveref}
\usepackage{enumitem}
\usepackage{xcolor}
\usepackage{booktabs}
\usepackage{fancyhdr}
\usepackage{longtable}

% Custom commands
\newcommand{\Id}{\mathbb{1}}
\renewcommand{\dd}{\mathrm{d}}
\newcommand{\ie}{\textit{i.e.}}
\newcommand{\eg}{\textit{e.g.}}
\newcommand{\Lam}{\Lambda}

% Header
\pagestyle{fancy}
\fancyhf{}
\fancyhead[L]{\small SCT Theory --- Literature Review}
\fancyhead[R]{\small MR-6: Curvature Expansion Convergence}
\fancyfoot[C]{\thepage}

\hypersetup{
  colorlinks=true,
  linkcolor=blue!60!black,
  citecolor=green!50!black,
  urlcolor=blue!70!black,
  pdfauthor={David Alfyorov},
  pdftitle={MR-6 Literature Review: Convergence of the Curvature Expansion},
  pdfsubject={Heat kernel, Seeley-DeWitt expansion, spectral action, convergence},
  pdfkeywords={heat kernel, spectral action, asymptotic expansion, Borel summability, NCG}
}

\title{MR-6: Literature Review\\
Convergence of the Curvature Expansion\\
in the Spectral Action\\[6pt]
\large Seeley--DeWitt Coefficients, Asymptoticity, and Borel Summability}
\author{David Alfyorov}
\date{March 2026}

\begin{document}
\maketitle

\begin{center}
\end{center}

\begin{abstract}
We review the literature on the convergence properties of the Seeley--DeWitt
(heat kernel) expansion of the spectral action, covering: (1)~computation and
growth of heat kernel coefficients by Gilkey, Avramidi, Vassilevich, and
others; (2)~exact spectral actions on symmetric spaces by Chamseddine--Connes;
(3)~asymptoticity and convergence studies by Estrada--Gracia-Bond\'{i}a--V\'{a}rilly,
Iochum--L\'{e}vy--Vassilevich, and Eckstein--Iochum; (4)~Borel summability
results by Harg\'{e} and Dunne; and (5)~the foundational spectral action
principle of Chamseddine--Connes and its relation to noncommutative geometry.
This review supports the MR-6 convergence analysis of Spectral Causal Theory.
\end{abstract}

\tableofcontents
\newpage


%======================================================================
\section{Introduction}\label{sec:intro}
%======================================================================

The spectral action principle~\cite{CC96} defines the bosonic action of
gravity and gauge fields through the spectrum of the Dirac operator.
Its connection to ordinary field theory is mediated by the heat kernel
expansion, whose convergence properties determine the validity of the
perturbative approach used in most applications.

This literature review surveys the principal results relevant to the
MR-6 consistency check of Spectral Causal Theory, organized by topic.
Section~\ref{sec:sdw-coeff} covers the computation of Seeley--DeWitt
coefficients.  Section~\ref{sec:exact-spectra} reviews exact spectral
computations on symmetric spaces.  Section~\ref{sec:convergence} addresses
convergence and asymptoticity.  Section~\ref{sec:borel} covers Borel
summability.  Section~\ref{sec:ncg-spectral-action} reviews the spectral
action in noncommutative geometry.  Section~\ref{sec:summary} summarizes
the findings relevant to MR-6.


%======================================================================
\section{Seeley--DeWitt Coefficient Computations}\label{sec:sdw-coeff}
%======================================================================

\subsection{Foundational work: Minakshisundaram, DeWitt, Gilkey}

The asymptotic expansion of the heat kernel trace was first studied by
Minakshisundaram and Pleijel~\cite{MP49}, who established the existence
of the expansion for the scalar Laplacian.  DeWitt~\cite{DeW65} extended
this to general Laplace-type operators on Riemannian manifolds, computing
the first three nonvanishing coefficients $a_0$, $a_2$, $a_4$.
Gilkey~\cite{G75,G95} placed the theory on a rigorous mathematical
footing, proving universality of the coefficients and computing $a_6$
in full generality (Theorem~4.8.16 of~\cite{G95}).

\subsection{Vassilevich: comprehensive review}

Vassilevich~\cite{V03} provides the standard reference for heat kernel
computations.  Key formulas (in the convention with $(4\pi)^{-d/2}$
prefactor absorbed):

\begin{itemize}
  \item $a_0 = (4\pi)^{-2}\int_M\dd^4x\,\sqrt{g}\,\mathrm{tr}(\Id)$
    \quad (eq.~4.1 of~\cite{V03})
  \item $a_2 = (4\pi)^{-2}\frac{1}{6}\int_M\dd^4x\,\sqrt{g}\,
    \mathrm{tr}(6E + R)$ \quad (eq.~4.3)
  \item $a_4$: the full 8-term expression involving $E^2$, $RE$,
    $\Omega^2$, $R^2$, $R_{\mu\nu}^2$, $R_{\mu\nu\rho\sigma}^2$
    (eq.~4.8)
  \item $a_6$: approximately 17--20 independent invariants (eq.~4.12)
\end{itemize}

Vassilevich also treats boundary contributions (Dirichlet, Robin,
mixed conditions) and reviews spectral geometry applications.

\subsection{Avramidi: covariant methods and growth estimates}

Avramidi~\cite{A00,A97} developed covariant techniques that yield
heat kernel coefficients algorithmically for arbitrary Laplace-type
operators.  A key result for MR-6 is the growth estimate
(Section~6 of~\cite{A00}):
\begin{equation}\label{eq:avramidi-growth}
  |a_{2k}(x,x)| \;\lesssim\; C(x)\,\Gamma\!\left(k + \frac{d}{2}\right)
  \left(\frac{R_{\max}}{\Lam_{\text{char}}}\right)^k,
\end{equation}
which confirms factorial (Gevrey-1) growth of the diagonal heat kernel
coefficients.

Avramidi~\cite{A08} also developed generating functions for heat kernel
coefficients on locally symmetric spaces, enabling exact spectral actions
on spaces such as $S^n$, $\mathbb{CP}^n$, and hyperbolic spaces.

\subsection{van de Ven: index-free computation of \texorpdfstring{$a_6$}{a6}}

Van de Ven~\cite{vdV98} computed the $a_6$ coefficient in index-free
notation, providing an independent check of Gilkey's result and a
computationally more efficient form.

\subsection{Branson--Gilkey: boundary contributions}

Branson and Gilkey~\cite{BG90} extended the computation of $a_6$ to
manifolds with boundary, including Dirichlet and Robin boundary
conditions.

\subsection{Salcedo: covariant derivative expansion}

Salcedo~\cite{S04} computed traced heat kernel coefficients with up to
6 covariant derivatives, providing diagonal coefficients with up to 4
derivatives in a form useful for effective action computations.

\subsection{Berline--Getzler--Vergne: mathematical foundations}

The monograph by Berline, Getzler, and Vergne~\cite{BGV04} provides
the rigorous mathematical framework for the heat kernel expansion via
parametrix construction.  Chapter~2 establishes the asymptotic nature
of the expansion and notes that the factorial growth of coefficients
precludes convergence on general manifolds.


%======================================================================
\section{Exact Spectra on Symmetric Spaces}\label{sec:exact-spectra}
%======================================================================

\subsection{Dirac spectrum on \texorpdfstring{$S^n$}{Sn}}

B\"{a}r~\cite{B96} computed the Dirac spectrum on space forms of
positive curvature.  For the round sphere $S^n$ of unit radius:
\begin{equation}
  \lambda_k = \pm\!\left(k + \frac{n}{2}\right), \quad
  m_k = 2^{\lfloor n/2\rfloor}\binom{n+k-1}{k}, \quad
  k = 0,1,2,\ldots
\end{equation}
Friedrich~\cite{F00} provides a comprehensive treatment of Dirac operators
on Riemannian manifolds, including spectral data for all compact symmetric
spaces.

For $S^4$ specifically: $\lambda_k = \pm(k+2)$ with $D^2$ multiplicity
$d_k = \frac{4}{3}(k+1)(k+2)(k+3)$, matching the formula
$P(m) = \frac{4}{3}m(m+1)(m-1)$ of Chamseddine--Connes~\cite{CC08}.

\subsection{Chamseddine--Connes: exact spectral action on \texorpdfstring{$S^3$}{S3}}

Chamseddine and Connes~\cite{CC08} computed the exact spectral action on
the round three-sphere $S^3$ and on $S^3 \times S^1$.  Key findings:

\begin{itemize}
  \item On $S^3$, only the first two Seeley--DeWitt terms (cosmological
    constant $+$ Einstein--Hilbert) are needed: $a_4$ and $a_6$ both vanish
    due to remarkable cancellations (conformally flat Einstein space).
  \item On $S^3 \times S^1$, the spectral action is given by two terms
    ``up to an astronomically small correction.''
  \item This ``uncanny precision'' is specific to maximally symmetric
    spaces and does \emph{not} generalize to arbitrary manifolds.
\end{itemize}

\subsection{Chamseddine--Connes: Robertson--Walker computation}

Chamseddine and Connes~\cite{CC11} computed Seeley--DeWitt coefficients
$a_0$ through $a_{10}$ for general Robertson--Walker metrics using a
combination of Euler--Maclaurin summation and Feynman--Kac methods.
The Bernoulli number structure in the sphere coefficients confirms
factorial growth of the expansion.

\subsection{Fathizadeh--Ghorbanpour--Khalkhali: rationality}

Fathizadeh, Ghorbanpour, and Khalkhali~\cite{FGK14} proved the
Chamseddine--Connes rationality conjecture: all coefficients in
$a_{2n}$ for Robertson--Walker metrics are rational numbers.  They
extended computations up to $a_{12}$ using pseudodifferential calculus
and confirmed agreement with~\cite{CC11} up to $a_{10}$.

\subsection{Fathizadeh--Marcolli: arithmetic structure}

Fathizadeh and Marcolli~\cite{FM16} showed that the coefficients of
the spectral action for Robertson--Walker spacetimes can be interpreted
as periods of mixed Tate motives, revealing deep arithmetic structure
in the heat kernel expansion.

\subsection{Other symmetric spaces}

The scalar Laplacian spectrum on $S^4$ is $\mu_l = l(l+3)$ with
multiplicity $d_l = (2l+3)(l+1)(l+2)/6$~\cite{BGM71}.  The Dirac
spectrum on $\mathbb{CP}^2$ was computed by Seifarth and
Semmelmann~\cite{SS93}.  On the flat torus $T^4$, the Dirac spectrum
reduces to lattice sums analyzable via Jacobi theta functions.


%======================================================================
\section{Convergence and Asymptoticity}\label{sec:convergence}
%======================================================================

\subsection{General consensus}

The heat kernel expansion is \textbf{asymptotic, not convergent}, on
general Riemannian manifolds.  This consensus is established across all
major references:

\begin{itemize}
  \item Vassilevich~\cite{V03} uses asymptotic equivalence (``$\sim$'')
    throughout; Section~8 discusses resummation approaches precisely
    because the series diverges.
  \item Avramidi~\cite{A00} emphasizes the expansion is valid only as
    $t \to 0^+$, with factorial growth precluding convergence.
  \item Gilkey~\cite{G95} states the expansion is asymptotic in the
    sense of Poincar\'{e}: for any $N$, the remainder after $N$ terms
    is $O(t^{(N+1-d)/2})$ as $t \to 0^+$.
  \item Berline--Getzler--Vergne~\cite{BGV04} prove the asymptotic
    expansion via parametrix construction; the factorial growth precludes
    convergence.
\end{itemize}

\subsection{Estrada--Gracia-Bond\'{i}a--V\'{a}rilly (1998)}

Estrada, Gracia-Bond\'{i}a, and V\'{a}rilly~\cite{EGV98} provided a
rigorous mathematical framework for the spectral action as a
\textbf{C\'{e}saro asymptotic development}:

\begin{itemize}
  \item The C\'{e}saro and parametric behaviours of distributions at
    infinity are shown to be equivalent modulo the moment asymptotic
    expansion.
  \item The bosonic action functional of Chamseddine--Connes is validated
    as a C\'{e}saro asymptotic development for $\Lam \to \infty$.
  \item This provides the rigorous justification for using the asymptotic
    expansion, but confirms it is \emph{not} a convergent series.
\end{itemize}

\subsection{Iochum--L\'{e}vy--Vassilevich (2011): beyond weak-field}

Iochum, L\'{e}vy, and Vassilevich~\cite{ILV11} went beyond the heat
kernel expansion to compute the \emph{exact} spectral action to second
order in the field strength.  Key results:

\begin{enumerate}
  \item The exact spectral action decays as $\sim 1/p^4$ in the UV
    (Theorem~3.4), in contrast to the polynomial growth suggested by
    the truncated expansion.
  \item ``The theory is not super-renormalizable in contrast to'' the
    truncated expansion --- a qualitative difference between exact and
    perturbative results.
  \item The coefficient $\Lam^4 c_f$ controlling UV asymptotics is
    identical to the one multiplying $a_0$, revealing an ``intimate
    relation between leading terms in the low and high momenta
    asymptotics.''
\end{enumerate}

This demonstrates that the asymptotic expansion can be qualitatively
misleading for finite $\Lam$.

\subsection{Iochum--L\'{e}vy--Vassilevich (2012): convergence criteria}

In a companion paper, Iochum, L\'{e}vy, and Vassilevich~\cite{ILV12}
directly studied convergence criteria for spectral action expansions,
computing the heat kernel on the torus to second order in the gauge
connection and deriving explicit convergence conditions.

\subsection{van Nuland--van Suijlekom (2021): absolute convergence}

Van Nuland and van Suijlekom~\cite{vNvS21} showed that the spectral
action, perturbed by a gauge potential, can be written as a series of
Chern--Simons and Yang--Mills actions of all orders, with conditions
for absolute convergence.  While this concerns gauge perturbations
rather than the curvature expansion, the convergence methodology is
relevant.

\subsection{Hekkelman--McDonald--van Nuland (2024): rigorous asymptotics}

Hekkelman, McDonald, and van Nuland~\cite{HMvN24} established a
rigorous asymptotic expansion of
$\mathrm{Tr}\bigl(P\,e^{-t(D+V)^2}\bigr)$ via multiple operator integral
techniques, providing the most modern mathematical framework for
spectral action asymptotics.

\subsection{Mistry--Pinzul--Rachwal (2020): higher derivatives}

Mistry, Pinzul, and Rachwal~\cite{MPR20} computed the complete spectral
action up to 6-derivative terms, demonstrating the rigid structure of
higher-derivative gravity arising from the spectral action.

\subsection{Eckstein--Iochum (2019): comprehensive treatment}

The monograph by Eckstein and Iochum~\cite{EI19} provides the most
comprehensive treatment of the spectral action, including:

\begin{itemize}
  \item Chapter~3: ``Asymptotic Expansion of the Spectral Action''
    (rigorous framework)
  \item Careful distinction between convergent expansions (on specific
    spaces like $S^3$) and generally asymptotic expansions (on
    arbitrary manifolds)
  \item Sections on ``Convergent, Non-exact, Expansions'' (p.~84)
    and ``Exact Expansions'' (p.~85)
\end{itemize}

\subsection{van Suijlekom: textbook treatment}

Van Suijlekom~\cite{vS15} uses heat kernel methods to derive the
spectral action as ``an asymptotic expansion for the spectral action
on $M \times F$ in terms of local formulas (on $M$).''  The expansion
is explicitly described as asymptotic throughout.

\subsection{Summary of convergence status}

\begin{table}[h]
\centering
\begin{tabular}{lll}
\toprule
Space/Setting & Status & Key Reference \\
\midrule
General $M^4$ & Asymptotic (divergent) & \cite{V03,A00,BGV04,G95} \\
Round $S^3$ & Finite (terminates at 2 terms) & \cite{CC08} \\
$S^3 \times S^1$ & Finite $+$ exp.\ small & \cite{CC08} \\
Robertson--Walker & Asymptotic & \cite{CC11} \\
Flat torus $T^4$ & Borel summable & \cite{H13a,H13b} \\
Hyperbolic $H^{2n}$ & Resummable & \cite{D21} \\
Momentum space (exact) & $1/p^4$ UV decay & \cite{ILV11} \\
\bottomrule
\end{tabular}
\caption{Convergence status of the heat kernel expansion across
different settings.}
\label{tab:convergence-summary}
\end{table}


%======================================================================
\section{Borel Summability}\label{sec:borel}
%======================================================================

\subsection{Gevrey-1 class and the Borel transform}

The Seeley--DeWitt coefficients satisfy $|a_{2k}| \leq C^{k+1}(2k)!$,
which is Gevrey-1 class.  For such series, the Borel transform
$\mathcal{B}[\sum c_k t^k](\zeta) = \sum c_k \zeta^k / k!$ may converge
and the original function may be recoverable by Laplace transform of
the Borel sum.

\subsection{Harg\'{e} (2013): Borel summability on tori}

Harg\'{e}~\cite{H13a} proved Borel summability of the small-time
expansion of the heat kernel for the Schr\"{o}dinger operator
$-\Delta + V(x)$ on the torus $T^d$.  The mechanism uses the Poisson
representation of the heat kernel:
\begin{equation}
  K(t,x,y) = \sum_\gamma K_0(t,x,y+\gamma)\,e^{i\phi(\gamma)},
\end{equation}
where the sum runs over lattice vectors.  The Borel transform of the
asymptotic series equals this exact sum, establishing Borel summability.

Harg\'{e}~\cite{H13b} extended the result to first-order perturbations
$-\Delta + A_\mu\nabla^\mu$ of the Laplacian.  The complex extension
was studied in~\cite{H13c}.

\subsection{Dunne (2021): hyperbolic spaces}

Dunne~\cite{D21} studied the Borel properties of the heat kernel
expansion on even-dimensional hyperbolic spaces $H^{2n}$.  Key results:

\begin{itemize}
  \item The diagonal heat kernel expansion has alternating-sign factorially
    divergent coefficients with Bernoulli number structure.
  \item Resummation via incomplete gamma functions provides accurate
    extrapolation and analytic continuation (from $H^2$ to $S^2$ and from
    heat kernel to Schr\"{o}dinger kernel).
  \item The Borel properties differ between short-time and long-time
    expansions.
\end{itemize}

This demonstrates that Borel-like resummation is feasible on curved
symmetric spaces beyond the flat torus.

\subsection{General manifolds: open problem}

For general curved manifolds, Borel summability of the heat kernel
expansion has \textbf{not} been established.  Key obstructions include:

\begin{enumerate}
  \item The heat kernel on curved manifolds lacks a simple Poisson-type
    representation.
  \item The geodesic structure (Synge world function) creates non-trivial
    contributions.
  \item Singularities of the Borel transform from topology (instantons,
    non-perturbative effects) can obstruct summability.
\end{enumerate}

Avramidi~\cite{A00} notes that generating functions on symmetric spaces
suggest some summability structure, but no rigorous theorem exists for
general curved backgrounds.


%======================================================================
\section{The Spectral Action in Noncommutative Geometry}\label{sec:ncg-spectral-action}
%======================================================================

\subsection{Chamseddine--Connes: spectral action principle}

Chamseddine and Connes~\cite{CC96} introduced the spectral action principle:
the bosonic action is $\mathrm{Tr}(f(D^2/\Lam^2))$ where $D$ is the Dirac
operator of a spectral triple.  The heat kernel expansion yields the
Einstein--Hilbert action with cosmological constant, plus curvature-squared
corrections, plus the full Standard Model Lagrangian when the spectral
triple encodes the SM algebra~\cite{CC10}.

The Laplace representation
\begin{equation}
  S = \int_0^\infty \psi(t)\,\mathrm{Tr}(e^{-tD^2/\Lam^2})\,\dd t
\end{equation}
provides the non-perturbative definition when $f$ admits an inverse
Laplace transform $\psi$.

\subsection{The role of the cutoff function}

The cutoff function $f$ enters through its moments $f_n = \int_0^\infty
u^{n/2-1}f(u)\,\dd u$.  For physical applications:

\begin{itemize}
  \item The exponential cutoff $f(u) = e^{-u}$ gives $f_n = \Gamma(n/2)$.
  \item The moments $f_{4-2k}$ for $k \geq 3$ decay as $1/(k-2)!$, which
    can compensate the factorial growth of $a_{2k}$ on specific spaces.
  \item For generic cutoff functions with polynomial moments, no such
    compensation occurs, and the expansion genuinely diverges.
\end{itemize}

\subsection{Spectral action on noncommutative spaces}

The spectral action has been computed on several noncommutative spaces:

\begin{itemize}
  \item \textbf{Noncommutative torus:} Essouabri, Iochum, L\'{e}vy, and
    Sitarz~\cite{ELS12} computed the spectral action via zeta function
    methods, encountering Diophantine conditions.
  \item \textbf{Podle\'{s} sphere:} Eckstein, Iochum, and Sitarz~\cite{EIS15}
    computed the exact heat trace on the standard Podle\'{s} sphere.
  \item \textbf{Almost-commutative geometries:} The spectral action on
    $M \times F$ (where $F$ is a finite spectral triple encoding the Standard
    Model) yields the full gravitational $+$ SM action~\cite{vS15,CC10}.
\end{itemize}

\subsection{Relation to SCT form factors}

In Spectral Causal Theory, the one-loop spectral action produces nonlocal
form factors $F_1(\Box/\Lam^2)$, $F_2(\Box/\Lam^2,\xi)$ built from the
master function
$\varphi(z) = e^{-z/4}\sqrt{\pi/z}\,\mathrm{erfi}(\sqrt{z}/2)$.
The entireness of $\varphi(z)$ (NT-2) means the \emph{nonlocal} spectral
action is well-defined for all curvatures.  The curvature expansion is
obtained by Taylor-expanding these form factors around $z = 0$ --- this
localization step introduces the factorial divergence.


%======================================================================
\section{Summary and Implications for MR-6}\label{sec:summary}
%======================================================================

The literature establishes the following picture for the curvature expansion
of the spectral action:

\begin{enumerate}
  \item \textbf{The Seeley--DeWitt expansion is asymptotic (Gevrey-1).}
    This is the consensus of Gilkey, Vassilevich, Avramidi, and
    Berline--Getzler--Vergne.  The coefficients grow as $(2k)!$, precluding
    convergence on general manifolds.

  \item \textbf{Specific spaces exhibit better behaviour.}
    On $S^3$, the expansion terminates (Chamseddine--Connes).
    On the torus, it is Borel summable (Harg\'{e}).
    On hyperbolic spaces, it is resummable via incomplete gamma functions
    (Dunne).

  \item \textbf{The exact spectral action differs qualitatively from
    truncations.}  Iochum--L\'{e}vy--Vassilevich showed that the exact
    result decays as $1/p^4$ in momentum space, while the truncated
    expansion suggests polynomial growth --- a qualitative discrepancy.

  \item \textbf{The Laplace representation is the standard non-perturbative
    definition.}  It provides a well-defined spectral action for all
    curvatures, bypassing the divergent expansion.

  \item \textbf{Borel summability on general manifolds remains open.}
    Results exist for tori and hyperbolic spaces, but the general curved
    case is unsettled.

  \item \textbf{The SCT form factors are entire} (NT-2), providing a
    convergent representation in the momentum variable $z = \Box/\Lam^2$.
    The divergence arises only from the localization step.
\end{enumerate}

These results directly support the MR-6 analysis: the curvature expansion
is asymptotic but the underlying theory is well-defined through the
Laplace representation and entire form factors.


%======================================================================
\section*{Acknowledgements}
%======================================================================

workflow support throughout this project.


%======================================================================
% References
%======================================================================
\begin{thebibliography}{99}

% === Foundational: heat kernel coefficients ===
\bibitem{MP49}
S.~Minakshisundaram and {\AA}.~Pleijel,
\textit{Some properties of the eigenfunctions of the Laplace-operator
on Riemannian manifolds},
Canadian J.\ Math.\ \textbf{1} (1949) 242--256.

\bibitem{DeW65}
B.S.~DeWitt,
\textit{Dynamical Theory of Groups and Fields},
Gordon and Breach, 1965.

\bibitem{G75}
P.B.~Gilkey,
\textit{The spectral geometry of a Riemannian manifold},
J.\ Diff.\ Geom.\ \textbf{10} (1975) 601--618.

\bibitem{G95}
P.B.~Gilkey,
\textit{Invariance Theory, the Heat Equation, and the Atiyah--Singer
Index Theorem}, 2nd edition, CRC Press, 1995.

\bibitem{V03}
D.V.~Vassilevich,
\textit{Heat kernel expansion: user's manual},
Phys.\ Rept.\ \textbf{388} (2003) 279--360,
\href{https://arxiv.org/abs/hep-th/0306138}{arXiv:hep-th/0306138}.

\bibitem{A00}
I.G.~Avramidi,
\textit{Heat Kernel and Quantum Gravity},
Lecture Notes in Physics Monographs, vol.~64, Springer, 2000;
\href{https://arxiv.org/abs/hep-th/0008167}{arXiv:hep-th/0008167}.

\bibitem{A97}
I.G.~Avramidi,
\textit{Covariant techniques for computation of the heat kernel},
Rev.\ Math.\ Phys.\ \textbf{11} (1999) 947--980,
\href{https://arxiv.org/abs/hep-th/9704166}{arXiv:hep-th/9704166}.

\bibitem{A08}
I.G.~Avramidi,
\textit{Heat kernel on homogeneous bundles over symmetric spaces},
Commun.\ Math.\ Phys.\ \textbf{288} (2009) 963--1006,
\href{https://arxiv.org/abs/math/0701489}{arXiv:math/0701489}.

\bibitem{BGV04}
N.~Berline, E.~Getzler, and M.~Vergne,
\textit{Heat Kernels and Dirac Operators},
Springer, 2004 (corrected reprint).

\bibitem{vdV98}
A.E.M.~van de Ven,
\textit{Index-free heat kernel coefficients},
Class.\ Quantum Grav.\ \textbf{15} (1998) 2311--2344,
\href{https://arxiv.org/abs/hep-th/9708152}{arXiv:hep-th/9708152}.

\bibitem{BG90}
T.P.~Branson and P.B.~Gilkey,
\textit{The asymptotics of the Laplacian on a manifold with boundary},
Commun.\ Part.\ Diff.\ Eqns.\ \textbf{15} (1990) 245--272.

\bibitem{S04}
L.L.~Salcedo,
\textit{Covariant derivative expansion of the heat kernel},
Eur.\ Phys.\ J.~C \textbf{37} (2004) 511--523,
\href{https://arxiv.org/abs/hep-th/0409140}{arXiv:hep-th/0409140}.

% === Exact spectra on symmetric spaces ===
\bibitem{B96}
C.~B\"{a}r,
\textit{The Dirac operator on space forms of positive curvature},
J.\ Math.\ Soc.\ Japan \textbf{48} (1996) 69--83.

\bibitem{F00}
T.~Friedrich,
\textit{Dirac Operators in Riemannian Geometry},
Graduate Studies in Mathematics, vol.~25, AMS, 2000.

\bibitem{BGM71}
M.~Berger, P.~Gauduchon, and E.~Mazet,
\textit{Le spectre d'une vari\'{e}t\'{e} riemannienne},
Lecture Notes in Mathematics, vol.~194, Springer, 1971.

\bibitem{SS93}
H.~Seifarth and U.~Semmelmann,
\textit{The spectrum of the Dirac operator on
$\mathbb{CP}^2$},
preprint, 1993.

% === Exact spectral action computations ===
\bibitem{CC08}
A.H.~Chamseddine and A.~Connes,
\textit{The Uncanny Precision of the Spectral Action},
Commun.\ Math.\ Phys.\ \textbf{293} (2010) 867--897,
\href{https://arxiv.org/abs/0812.0165}{arXiv:0812.0165}.

\bibitem{CC11}
A.H.~Chamseddine and A.~Connes,
\textit{Spectral Action for Robertson--Walker metrics},
JHEP \textbf{10} (2012) 101,
\href{https://arxiv.org/abs/1105.4637}{arXiv:1105.4637}.

\bibitem{FGK14}
F.~Fathizadeh, A.~Ghorbanpour, and M.~Khalkhali,
\textit{Rationality of Spectral Action for Robertson--Walker Metrics},
JHEP \textbf{12} (2014) 064,
\href{https://arxiv.org/abs/1407.5972}{arXiv:1407.5972}.

\bibitem{FM16}
F.~Fathizadeh and M.~Marcolli,
\textit{Periods and motives in the spectral action of Robertson--Walker
spacetimes},
\href{https://arxiv.org/abs/1611.01815}{arXiv:1611.01815}, 2016.

% === Convergence and asymptoticity ===
\bibitem{EGV98}
R.~Estrada, J.M.~Gracia-Bond\'{i}a, and J.C.~V\'{a}rilly,
\textit{On summability of distributions and spectral geometry},
Commun.\ Math.\ Phys.\ \textbf{191} (1998) 219--248,
\href{https://arxiv.org/abs/funct-an/9702001}{arXiv:funct-an/9702001}.

\bibitem{ILV11}
B.~Iochum, C.~L\'{e}vy, and D.V.~Vassilevich,
\textit{Spectral action beyond the weak-field approximation},
Commun.\ Math.\ Phys.\ \textbf{316} (2012) 595--613,
\href{https://arxiv.org/abs/1108.3749}{arXiv:1108.3749}.

\bibitem{ILV12}
B.~Iochum, C.~L\'{e}vy, and D.V.~Vassilevich,
\textit{Global and local aspects of spectral actions},
J.\ Phys.\ A \textbf{45} (2012) 374020,
\href{https://arxiv.org/abs/1201.6637}{arXiv:1201.6637}.

\bibitem{vNvS21}
T.D.H.~van Nuland and W.D.~van Suijlekom,
\textit{Cyclic cocycles in the spectral action},
J.\ Noncommut.\ Geom.\ (2022),
\href{https://arxiv.org/abs/2104.09899}{arXiv:2104.09899}.

\bibitem{HMvN24}
E.-M.~Hekkelman, E.~McDonald, and T.D.H.~van Nuland,
\textit{Multiple operator integrals, pseudodifferential calculus,
and asymptotic expansions},
\href{https://arxiv.org/abs/2404.16338}{arXiv:2404.16338}, 2024.

\bibitem{MPR20}
R.~Mistry, A.~Pinzul, and L.~Rachwal,
\textit{Spectral action approach to higher derivative gravity},
Phys.\ Rev.\ D \textbf{102} (2020) 024033,
\href{https://arxiv.org/abs/2001.05975}{arXiv:2001.05975}.

\bibitem{EI19}
M.~Eckstein and B.~Iochum,
\textit{Spectral Action in Noncommutative Geometry},
SpringerBriefs in Mathematical Physics, vol.~27, 2019,
\href{https://arxiv.org/abs/1902.05306}{arXiv:1902.05306}.

\bibitem{vS15}
W.D.~van Suijlekom,
\textit{Noncommutative Geometry and Particle Physics},
Springer, 2015 (2nd edition 2024).

% === Borel summability ===
\bibitem{H13a}
T.~Harg\'{e},
\textit{Borel summation of the small time expansion of the heat kernel.
The scalar potential case},
\href{https://arxiv.org/abs/1301.7742}{arXiv:1301.7742}, 2013.

\bibitem{H13b}
T.~Harg\'{e},
\textit{Borel summation of the small time expansion of the heat kernel
with a vector potential},
\href{https://arxiv.org/abs/1302.0604}{arXiv:1302.0604}, 2013.

\bibitem{H13c}
T.~Harg\'{e},
\textit{Some remarks on the complex heat kernel in the scalar
potential case},
\href{https://arxiv.org/abs/1302.3512}{arXiv:1302.3512}, 2013.

\bibitem{D21}
G.V.~Dunne,
\textit{Borel Summation and Analytic Continuation of the Heat Kernel
on Hyperbolic Space},
\href{https://arxiv.org/abs/2109.03897}{arXiv:2109.03897}, 2021.

% === Spectral action principle ===
\bibitem{CC96}
A.H.~Chamseddine and A.~Connes,
\textit{The Spectral Action Principle},
Commun.\ Math.\ Phys.\ \textbf{186} (1997) 731--750,
\href{https://arxiv.org/abs/hep-th/9606001}{arXiv:hep-th/9606001}.

\bibitem{CC10}
A.H.~Chamseddine and A.~Connes,
\textit{Noncommutative geometry as a framework for unification},
Fortsch.\ Phys.\ \textbf{58} (2010) 553--600,
\href{https://arxiv.org/abs/1004.0464}{arXiv:1004.0464}.

% === Noncommutative spaces ===
\bibitem{EIS15}
M.~Eckstein, B.~Iochum, and A.~Sitarz,
\textit{Heat trace and spectral action on the standard Podle\'{s} sphere},
Commun.\ Math.\ Phys.\ \textbf{332} (2014) 627--668.

\bibitem{ELS12}
D.~Essouabri, B.~Iochum, C.~L\'{e}vy, and A.~Sitarz,
\textit{Spectral action on noncommutative torus},
J.\ Noncommut.\ Geom.\ \textbf{2} (2008) 53--123.

\end{thebibliography}

\end{document}
